\documentclass[a4paper,12pt]{article}
\usepackage[utf8]{inputenc}
\usepackage[T1]{fontenc}
\usepackage[english]{babel}
\usepackage{geometry}
\usepackage{fancyhdr}
\usepackage{tcolorbox}
\usepackage{lmodern}
\usepackage{lastpage}
\usepackage{datetime}
\usepackage{titlesec}
\usepackage{changepage}
\usepackage[hidelinks]{hyperref}
\usepackage{graphicx}
\usepackage{listings}
\usepackage{url}
\usepackage{breakurl}
\usepackage{seqsplit}
\usepackage{adjustbox}
\usepackage{array}
\usepackage{longtable}
\usepackage{booktabs}
\usepackage{listingsutf8}
\lstset{inputencoding=utf8/latin2}

\geometry{margin=2.5cm}
\pagestyle{fancy}
\fancyhf{}
\fancyhead[L]{NFC SERVICES}
\fancyhead[R]{2025-07-08}
\fancyfoot[R]{\textbf{Page \thepage\ of \pageref{LastPage}}}

% Configuration for long code lines
\lstset{
    breaklines=true,
    breakatwhitespace=true,
    basicstyle=\footnotesize\ttfamily,
    frame=single,
    numbers=left,
    numberstyle=\tiny,
    stepnumber=1,
    numbersep=5pt,
    tabsize=2,
    showstringspaces=false,
    showtabs=false,
    captionpos=b,
    xleftmargin=2em,
    framexleftmargin=1.5em
}

\begin{document}
\vspace*{1cm}

\begin{center}
\begin{tcolorbox}[width=0.8\textwidth, sharp corners, boxrule=0.5pt]
    \centering
    \textbf{\Large Technical Documentation \\ NFC Services} \\[1.5em]
    \textbf{Version:} 1.0 \\
    \textbf{Release:} A \\[1.5em]
    \textbf{Author:} Michal Itrych \\[1.5em]
    \textbf{Publication Date:} 08-07-2025
\end{tcolorbox}
\end{center}

\newpage
\tableofcontents
\newpage

\section{Introduction and System Overview}

\subsection{Purpose of Documentation}
This technical documentation describes the implementation of the NFC Services system for Flutter applications, enabling communication with MIFARE DESFire and MIFARE Plus SL3 cards.

\subsection{System Architecture}
The NFC Services system consists of the following main components:

\begin{itemize}
    \item \textbf{DesfireNfcServices} -- class handling operations on MIFARE DESFire cards
    \item \textbf{PlusNfcServices} -- class handling operations on MIFARE Plus SL3 cards
    \item \textbf{Shared utilities} -- tools for cryptography, NFC and conversion (ASCII, hex, bin).
\end{itemize}

\subsection{Supported Card Types}
The system supports the following NFC card types:
\begin{itemize}
    \item \textbf{MIFARE DESFire} -- cards with advanced cryptographic features
    \item \textbf{MIFARE Plus SL3} -- cards with additional security mechanisms
\end{itemize}

\newpage

\subsection{Directory Structure}
\begin{lstlisting}[language=bash]
lib/
  nfc_services/
    desfire_services.dart
    plus_services.dart
    shared/
      crypto_utilities.dart
      nfc_utilities.dart
      binary_conversion.dart
  example_use/
    desfire_screen.dart
    plus_screen.dart
\end{lstlisting}

\textbf{File descriptions:}
\begin{itemize}
    \item \texttt{desfire\_services.dart} -- DESFire card handling
    \item \texttt{plus\_services.dart} -- MIFARE Plus card handling
    \item \texttt{crypto\_utilities.dart} -- Cryptographic tools
    \item \texttt{nfc\_utilities.dart} -- Basic NFC operations
    \item \texttt{binary\_conversion.dart} -- Binary conversions
    \item \texttt{desfire\_screen.dart} -- DESFire usage example
    \item \texttt{plus\_screen.dart} -- MIFARE Plus usage example
\end{itemize}

\subsection{External Dependencies}
The system uses the following libraries:

\begin{itemize}
    \item \texttt{flutter\_nfc\_kit} – for basic NFC communication,
    \item \texttt{convert} – for data conversion (e.g., between hex, binary and ASCII formats),
    \item \texttt{encrypt} – for basic cryptographic operations,
    \item \texttt{pointycastle} – for advanced cryptographic functions.
\end{itemize}

To use these libraries, add them to the \texttt{pubspec.yaml} file:

\begin{lstlisting}[language=bash]
dependencies:
  flutter_nfc_kit: ^3.3.1
  convert: ^3.0.1
  encrypt: ^5.0.1
  pointycastle: ^3.9.1
\end{lstlisting}

\subsection{Supported Platforms}
The system is compatible with:
\begin{itemize}
    \item \textbf{Android} -- SDK 26+ (Android 8.0+)
    \item \textbf{iOS} -- iOS 14.6+
\end{itemize}

\newpage

\section{Installation and Configuration}

\subsection{System Requirements}

\subsubsection{Hardware Requirements}
\begin{itemize}
    \item \textbf{NFC-enabled device} -- phone or tablet with built-in NFC module
    \item \textbf{Android} -- device with Android 8.0 (API 26) or newer
    \item \textbf{iOS} -- iPhone with iOS 14.6 or newer (iPhone 7 and newer)
\end{itemize}

\subsubsection{Software Requirements}
\begin{itemize}
    \item \textbf{Flutter SDK} -- version 3.7.2 or newer
    \item \textbf{Dart SDK} -- version 3.0.0 or newer
    \item \textbf{Android Studio} -- version 2022.1 or newer (for Android development)
    \item \textbf{Xcode} -- version 14.0 or newer (for iOS development)
    \item \textbf{Gradle} -- version 7.0 or newer (for Android)
\end{itemize}

\subsubsection{NFC Requirements}
\begin{itemize}
    \item \textbf{NFC Standard} -- device must support ISO/IEC 14443 standard
    \item \textbf{Reader Mode} -- device must support NFC reader mode
\end{itemize}

\subsection{Package Installation}

\subsubsection{Adding Dependencies}
To install NFC Services in your Flutter project, add the following dependencies to your \texttt{pubspec.yaml} file:

\begin{lstlisting}
dependencies:
  flutter:
    sdk: flutter
  
  # NFC Services dependencies
  flutter_nfc_kit: ^3.3.1
  convert: ^3.0.1
  encrypt: ^5.0.1
  pointycastle: ^3.9.1
  
  # Additional dependencies for your app
  cupertino_icons: ^1.0.2
\end{lstlisting}

\subsubsection{Installing Dependencies}
After adding the dependencies, run the following commands in the terminal:

\begin{lstlisting}[language=bash]
flutter pub get

flutter doctor

flutter run
\end{lstlisting}

\subsection{Permission Configuration}

\subsubsection{iOS Configuration}
Add the following keys to the \texttt{ios/Runner/Info.plist} file:

\begin{lstlisting}[language=xml]
<?xml version="1.0" encoding="UTF-8"?>
<!DOCTYPE plist PUBLIC "-//Apple//DTD PLIST 1.0//EN" "http://www.apple.com/DTDs/PropertyList-1.0.dtd">
<plist version="1.0">
<dict>
    <!-- Other keys... -->
    
    <key>NFCReaderUsageDescription</key>
	<string>NFC capabilities to function properly.</string>
</dict>
</plist>
\end{lstlisting}

Also add the following keys to the \texttt{ios/Runner/Runner.entitlements} file:

\begin{lstlisting}[language=xml]
<?xml version="1.0" encoding="UTF-8"?>
<!DOCTYPE plist PUBLIC "-//Apple//DTD PLIST 1.0//EN" "http://www.apple.com/DTDs/PropertyList-1.0.dtd">
<plist version="1.0">
<dict>
	<key>com.apple.developer.nfc.readersession.formats</key>
	<array>
		<string>TAG</string>
	</array>
</dict>
</plist>
\end{lstlisting}

\newpage

\subsection{Installation Verification}

\subsubsection{Checking NFC Availability}
To verify that NFC is properly configured, use the following code:

\begin{lstlisting}
import 'package:flutter_nfc_kit/flutter_nfc_kit.dart';

Future<void> checkNfcAvailability() async {
  try {
    final availability = await FlutterNfcKit.nfcAvailability;
    
    switch (availability) {
      case NFCAvailability.available:
        print('NFC is available and ready to use');
        break;
      case NFCAvailability.disabled:
        print('NFC is disabled - please enable in device settings');
        break;
      case NFCAvailability.not_supported:
        print('NFC is not supported on this device');
        break;
    }
  } catch (e) {
    print('Error checking NFC availability: $e');
  }
}
\end{lstlisting}

\newpage

\subsubsection{Basic Functionality Test}
Create a simple test to verify everything is working:

\begin{lstlisting}
import 'package:flutter/material.dart';
import 'nfc_services/desfire_services.dart';

class NfcTestScreen extends StatefulWidget {
  @override
  _NfcTestScreenState createState() => _NfcTestScreenState();
}

class _NfcTestScreenState extends State<NfcTestScreen> {
  String _testResult = '';

  Future<void> _testNfcConnection() async {
    try {
      final desfire = DesfireNfcServices(debugMode: true);
      
      setState(() {
        _testResult = 'Testing NFC connection...';
      });
      
      // Test basic NFC availability
      final result = await desfire.readDesfire(
        'F40001', // Application ID
        '00',     // File ID
        '00',     // Key Number
        '00000000000000000000000000000000', // Auth Key
      );
      
      setState(() {
        if (result.isSuccess) {
          _testResult = 'NFC test successful! Data: ${result.data}';
        } else {
          _testResult = 'NFC test failed: ${result.errorCode?.description}';
        }
      });
    } catch (e) {
      setState(() {
        _testResult = 'Error during NFC test: $e';
      });
    }
  }

  @override
  Widget build(BuildContext context) {
    return Scaffold(
      appBar: AppBar(title: Text('NFC Test')),
      body: Padding(
        padding: EdgeInsets.all(16.0),
        child: Column(
          children: [
            ElevatedButton(
              onPressed: _testNfcConnection,
              child: Text('Test NFC Connection'),
            ),
            SizedBox(height: 16),
            Text(_testResult),
          ],
        ),
      ),
    );
  }
}
\end{lstlisting}

\subsection{Installation Troubleshooting}

\subsubsection{Common Android Issues}
\begin{itemize}
    \item \textbf{Issue:} "NFC not available" on NFC-enabled device
    \begin{itemize}
        \item \textbf{Solution:} Check if NFC is enabled in device settings
        \item \textbf{Solution:} Ensure the app has NFC permissions
    \end{itemize}
    
    \item \textbf{Issue:} "Permission denied" during NFC operations
    \begin{itemize}
        \item \textbf{Solution:} Check if permissions are properly added to AndroidManifest.xml
        \item \textbf{Solution:} Restart the app after adding permissions
    \end{itemize}
    
    \item \textbf{Issue:} "Gradle build failed"
    \begin{itemize}
        \item \textbf{Solution:} Check if all dependencies are properly added
        \item \textbf{Solution:} Run \texttt{flutter clean} and \texttt{flutter pub get}
    \end{itemize}
\end{itemize}

\subsubsection{Common iOS Issues}
\begin{itemize}
    \item \textbf{Issue:} "NFC not supported" on iPhone
    \begin{itemize}
        \item \textbf{Solution:} Ensure you're using iPhone 7 or newer
        \item \textbf{Solution:} Check if iOS version is 14.6+
    \end{itemize}
    
    \item \textbf{Issue:} "NFCReaderUsageDescription missing"
    \begin{itemize}
        \item \textbf{Solution:} Add NFC usage description to Info.plist
        \item \textbf{Solution:} Ensure the key is properly formatted
    \end{itemize}
    
    \item \textbf{Issue:} "Entitlements error"
    \begin{itemize}
        \item \textbf{Solution:} Check if Runner.entitlements file is properly configured
        \item \textbf{Solution:} Ensure you have proper developer certificates
    \end{itemize}
\end{itemize}

\subsubsection{Dependency Issues}
\begin{itemize}
    \item \textbf{Issue:} Library version conflicts
    \begin{itemize}
        \item \textbf{Solution:} Use \texttt{flutter pub deps} to check dependencies
        \item \textbf{Solution:} Update Flutter to the latest version
    \end{itemize}
    
    \item \textbf{Issue:} "Package not found"
    \begin{itemize}
        \item \textbf{Solution:} Check if pubspec.yaml is properly formatted
        \item \textbf{Solution:} Run \texttt{flutter pub cache repair}
    \end{itemize}
\end{itemize}

\section{Architecture and Main Components}

\subsection{Architecture Overview}
The NFC Services system has a modular structure – it is divided into a card communication layer, a cryptographic tools layer, and a presentation layer (e.g., example screens).

\subsection{Main Components}

\begin{itemize}
    \item \textbf{DesfireNfcServices} – class responsible for handling MIFARE DESFire cards, including authentication, file reading, error handling, and logging.
    
    \item \textbf{PlusNfcServices} – class responsible for handling MIFARE Plus SL3 cards (works similarly to \texttt{DesfireNfcServices}, but is based on a different protocol).
    
    \item \textbf{Shared tools (\texttt{shared/})} – set of tools used by both services mentioned above:
    \begin{itemize}
        \item \texttt{crypto\_utilities.dart} – cryptographic functions, including AES encryption, random number generation, CMAC calculation, and byte operations,
        \item \texttt{nfc\_utilities.dart} – functions supporting NFC session handling, such as availability checking, card detection, command sending, and error handling,
        \item \texttt{binary\_conversion.dart} – helper functions for binary and hexadecimal (hex) data conversion.
    \end{itemize}
    
\end{itemize}

\subsection{Dependency Schema}
\begin{itemize}
    \item \texttt{desfire\_services.dart} and \texttt{plus\_services.dart} use tools from the \texttt{shared/} directory.
    \item All NFC operations are based on the \texttt{flutter\_nfc\_kit} library.
\end{itemize}

\section{DesfireNfcServices Module}

\subsection{Class Description}
\texttt{DesfireNfcServices} is the main class responsible for handling MIFARE DESFire cards. It implements the complete communication cycle: from checking NFC availability, through application selection and authentication, to data reading, error handling, and log generation.

\subsection{Result Handling System}
Each asynchronous method returns a result as an object derived from the \texttt{DesfireResult<T>} class

\begin{itemize}
    \item \texttt{DesfireSuccess<T>} -- successful operation:
    \begin{itemize}
        \item \texttt{data} -- operation result of type \texttt{T}
        \item \texttt{isSuccess} -- always \texttt{true}
        \item \texttt{errorCode} -- always \texttt{null}
        \item \texttt{message} -- always \texttt{null}
        \item \texttt{exception} -- always \texttt{null}
    \end{itemize}
    \item \texttt{DesfireError<T>} -- operation error:
    \begin{itemize}
        \item \texttt{data} -- always \texttt{null}
        \item \texttt{isSuccess} -- always \texttt{false}
        \item \texttt{errorCode} -- error code of type \texttt{DesfireErrorCode}
        \item \texttt{message} -- optional error description (\texttt{String?})
        \item \texttt{exception} -- optional exception (\texttt{Exception?})
    \end{itemize}
\end{itemize}

\textbf{Common properties:}
\begin{itemize}
    \item \texttt{isSuccess} -- \texttt{bool} determining whether the operation was successful
    \item \texttt{data} -- operation result or \texttt{null} in case of error
    \item \texttt{errorCode} -- error code or \texttt{null} in case of success
    \item \texttt{message} -- error description or \texttt{null} in case of success
    \item \texttt{exception} -- exception or \texttt{null} in case of success
\end{itemize}

\newpage

\subsection{Types of Operation Results}
The NFC Services system uses the \textbf{result handling pattern} for consistent management of operation successes and errors:

\subsubsection*{Result Structure}
\begin{lstlisting}[language=Java]
// Success
DesfireSuccess<String>(
  data: "48656C6C6F20576F726C642100000000"
)

// Error
DesfireError<String>(
  error: DesfireErrorCode.authenticationFailed,
  message: "Authentication failed - check your key"
)
\end{lstlisting}

\subsubsection*{Checking Result Type}
\begin{lstlisting}[language=Java]
final result = await desfire.readDesfire(...);

// Method 1:
if (result.isSuccess) {
    final data = result.data!; // Data is available
    print('Read data: $data');
} else {
    final error = result.errorCode!; // Error code is available
    print('Error: ${error.description}');
}

// Method 2:
switch (result) {
    case DesfireSuccess(data: final data):
        print('Read data: $data');
    case DesfireError(error: final error, message: final message):
        print('Error: ${error.description}');
        if (message != null) print('Details: $message');
}

// Method 3:
if (!result.isSuccess) {
    final error = result.errorCode!; // Error code is available
    print('Error: ${error.description}');
} else {
    final data = result.data!; // Data is available
    print('Read data: $data');
}
\end{lstlisting}

\newpage

\subsubsection*{Properties Available for All Results}
\begin{itemize}
    \item \texttt{isSuccess} -- \texttt{bool} - whether the operation succeeded
    \item \texttt{data} -- \texttt{T?} - result data (null for errors)
    \item \texttt{errorCode} -- \texttt{DesfireErrorCode?} - error code (null for successes)
\end{itemize}

\subsection{Constructor}
\begin{lstlisting}[language=Java]
DesfireNfcServices({
  bool debugMode = false, 
  int timeout = 7, 
  String alertMessage, 
  String multipleTagMessage
})
\end{lstlisting}

\textbf{Parameters:}
    \begin{itemize}
        \item \texttt{debugMode} (\texttt{bool}) -- enables log generation
        \item \texttt{timeout} (\texttt{int}) -- NFC operation timeout (seconds)
        \item \texttt{alertMessage} (\texttt{String}) -- message during scanning (iOS only)
        \item \texttt{multipleTagMessage} (\texttt{String}) -- message for multiple tags (iOS only)
\end{itemize}

\subsection{Public Methods}
\begin{lstlisting}[language=Java]
Future<DesfireResult<String>> readDesfire(
  String applicationId, 
  String fileId, 
  String keyNumber, 
  String authKey
)
\end{lstlisting}

\textbf{Parameters:}
    \begin{itemize}
        \item \texttt{applicationId} (\texttt{String}) -- application identifier (6 hex characters)
        \item \texttt{fileId} (\texttt{String}) -- file identifier (0-31)
        \item \texttt{keyNumber} (\texttt{String}) -- key number (0-13)
        \item \texttt{authKey} (\texttt{String}) -- authentication key (32 hex characters)
    \end{itemize}

\subsection{Helper Types}
\begin{itemize}
    \item \texttt{DesfireAuthResult} -- authentication result (codes, random numbers, session key)
    \item \texttt{DesfireFileSettings} -- file settings (type, communication mode, access rights)
    \item \texttt{DesfireFirstAuthResult} -- first authentication step result (RandomB)
\end{itemize}

\subsection{Private Methods}

\subsubsection*{Validation and Error Handling}
\begin{itemize}
    \item \texttt{DesfireErrorCode \_validateDesfireReadParameters(...)}\\
    Verifies the correctness of input parameters (e.g., length, hex format)
    
    \item \texttt{String \_getDesfireErrorMessage(String code)}\\
    Maps error code to readable text description
    
    \item \texttt{void \_validateCardResponse(...)}\\
    Checks card response correctness (e.g., length, status)
    
    \item \texttt{void \_log(String message, \{String level = 'INFO'\})}\\
    Logs messages (active only in \texttt{debugMode})
\end{itemize}

\subsubsection*{Card Operations}
\begin{itemize}
    \item \texttt{Future<DesfireResult<void>> \_selectApplication(String applicationId)}\\
    Sends application selection command to the card
    
    \item \texttt{Future<DesfireResult<DesfireFileSettings>> \_getFileSettings(String fileId)}\\
    Gets file settings (type, access rights, communication mode)
    
    \item \texttt{Future<DesfireResult<String>> \_readEncryptedData(...)}\\
    Reads and decrypts data from file
    
    \item \texttt{Future<String> \_decryptFileData(...)}\\
    Decrypts data block from file based on session key and CMAC
\end{itemize}

\subsubsection*{Authentication}
\begin{itemize}
    \item \texttt{Future<DesfireResult<DesfireFirstAuthResult>> \_processFirstAuthStep(...)}\\
    First authentication step (getting and decrypting RandomB)
    
    \item \texttt{Future<DesfireResult<DesfireAuthResult>> \_processSecondAuthStep(...)}\\
    Second authentication step (RandomA, session, verification)
    
    \item \texttt{Future<DesfireResult<DesfireAuthResult>> \_authenticateWithApplication(...)}\\
    Complete authentication with retry handling
    
    \item \texttt{bool \_verifyRotatedRandomA(Uint8List received, Uint8List expected)}\\
    Compares RandomA values in a timing-attack resistant manner
    
    \item \texttt{Uint8List \_generateSessionKey(Uint8List randomA, Uint8List randomB)}\\
    Generates session key based on RandomA and RandomB
    
    \item \texttt{String \_to3ByteLEHex(int value)}\\
    Converts number to 3-byte little-endian hex (for command parameters)
\end{itemize}

\newpage

\subsection{Error Handling}
Each method returns a detailed error code (\texttt{DesfireErrorCode}) and optional description and exception.

\begin{table}[h!]
\centering
\begin{tabular}{|l|l|p{6cm}|}
\hline
\textbf{Error Code} & \textbf{Code} & \textbf{Description} \\
\hline
\texttt{success} & 00 & Operation completed successfully \\
\hline
\texttt{invalidApplicationId} & DAID & Invalid application identifier (must be 6 hex characters) \\
\hline
\texttt{invalidFileId} & DFID & Invalid file identifier \\
\hline
\texttt{invalidKeyNumber} & DKN & Invalid key number \\
\hline
\texttt{invalidAuthKey} & DAK & Invalid authentication key (must be 32 hex characters) \\
\hline
\texttt{timeoutError} & DT & Operation timeout exceeded \\
\hline
\texttt{nfcUnavailable} & DN & NFC not available \\
\hline
\texttt{sessionError} & DSN & NFC session error \\
\hline
\texttt{unexpectedState} & DU & Unexpected application state \\
\hline
\texttt{userCancelled} & DUC & User cancelled operation \\
\hline
\texttt{cardError} & CARD & Card error \\
\hline
\texttt{authenticationFailed} & DAF & Authentication error \\
\hline
\texttt{sessionKeyError} & DSE & Session key generation error \\
\hline
\texttt{cmacError} & DCM & CMAC calculation error \\
\hline
\texttt{dataDecryptionError} & DDD & Data decryption error \\
\hline
\texttt{invalidResponse} & DIR & Invalid card response \\
\hline
\texttt{cardConnectionLost} & DCCL & Card connection lost - keep card steady \\
\hline
\texttt{communicationError} & DCE & Card communication error \\
\hline
\texttt{applicationNotFound} & DANF & Application not found \\
\hline
\texttt{fileNotFound} & DFNF & File not found \\
\hline
\texttt{permissionDenied} & DPD & Access denied \\
\hline
\texttt{conditionsNotSatisfied} & DCNS & Usage conditions not satisfied - command code not available in current card state \\
\hline
\texttt{unknownError} & UE & Unknown error. Try again \\
\hline
\end{tabular}
\caption{\texttt{DesfireErrorCode} error codes and their meanings}
\end{table}

\newpage

\subsection{Usage Example}
\begin{lstlisting}[language=Java]
final desfire = DesfireNfcServices(debugMode: true);
final result = await desfire.readDesfire(
  'A1B2C3',
  '01',
  '00',
  '0123456789ABCDEF0123456789ABCDEF',
);

if (result.isSuccess) {
  print('Block data: ${result.data}');
  print('Operation successful: ${result.isSuccess}');
} else {
  print('Error code: ${result.errorCode}');
  print('Error message: ${result.message}');
  if (result.exception != null) {
    print('Exception: ${result.exception}');
  }
}
\end{lstlisting}

\subsection{What readDesfire Can Return}

The \texttt{readDesfire} method returns an object of type \texttt{DesfireResult<String>}, which can be one of two types:

\subsubsection*{DesfireSuccess<String> - Successful Operation}
\begin{itemize}
    \item \texttt{isSuccess} = \texttt{true}
    \item \texttt{data} = \texttt{String} - data read from file in hex format
    \item \texttt{errorCode} = \texttt{null}
    \item \texttt{message} = \texttt{null}
    \item \texttt{exception} = \texttt{null}
\end{itemize}

\textbf{Success example:}
\begin{lstlisting}[language=Java]
DesfireSuccess<String>(
  data: "48656C6C6F20576F726C642100000000" // "Hello World!" in hex
)
\end{lstlisting}

\newpage

\subsubsection*{DesfireError<String> - Operation Error}
\begin{itemize}
    \item \texttt{isSuccess} = \texttt{false}
    \item \texttt{data} = \texttt{null}
    \item \texttt{errorCode} = \texttt{DesfireErrorCode} - error code
    \item \texttt{message} = \texttt{String?} - optional error description
    \item \texttt{exception} = \texttt{Object?} - optional exception
\end{itemize}

\textbf{Error examples:}
\begin{lstlisting}[language=Java]
// Error: parameter validation
DesfireError<String>(
  error: DesfireErrorCode.invalidApplicationId,
  message: "Invalid application ID (must be 6 hex characters)"
)

// Error: NFC not available
DesfireError<String>(
  error: DesfireErrorCode.nfcUnavailable,
  message: "Please enable NFC in your device settings"
)

// Error: authentication failed
DesfireError<String>(
  error: DesfireErrorCode.authenticationFailed,
  message: "Authentication failed - check your key"
)

// Error: connection lost
DesfireError<String>(
  error: DesfireErrorCode.cardConnectionLost,
  message: "Card connection lost - please keep card steady"
)

// Error: unknown error with exception
DesfireError<String>(
  error: DesfireErrorCode.unknownError,
  message: "Unknown error. Try again",
  exception: TimeoutException("Operation timed out")
)
\end{lstlisting}

\newpage

\subsubsection*{Returned Data Format}
In case of success, \texttt{data} contains:
\begin{itemize}
    \item \textbf{Format:} String in hex format (32 characters = 16 bytes)
    \item \textbf{Example:} \texttt{"48656C6C6F20576F726C642100000000"}
    \item \textbf{Decryption:} Data is automatically decrypted before returning
    \item \textbf{Encoding:} Data is in little-endian format according to DESFire standard
\end{itemize}

\subsubsection*{Error Handling in Application}
\begin{lstlisting}[language=Java]
final result = await desfire.readDesfire(appId, fileId, keyNum, authKey);

if (result.isSuccess) {
    // Success
    final blockData = result.data!;
    print('Read data: $blockData');
} else {
    // Error
    final errorCode = result.errorCode!;
    final errorMessage = result.message;
    
    switch (errorCode) {
        case DesfireErrorCode.nfcUnavailable:
            showDialog('Enable NFC in device settings');
            break;
        case DesfireErrorCode.cardConnectionLost:
            showDialog('Keep the card steady near the phone');
            break;
        case DesfireErrorCode.authenticationFailed:
            showDialog('Invalid authentication key');
            break;
        default:
            showDialog('Error: ${errorMessage ?? errorCode.description}');
    }
}
\end{lstlisting}

\section{PlusNfcServices Module}

\subsection{Class Description}
\texttt{PlusServices} is the main class responsible for handling MIFARE Plus SL3 cards. This class implements the complete authentication protocol and data reading from MIFARE Plus cards, using AES128 cryptographic mechanisms.

\subsection{Result Handling System}
Each asynchronous method returns a result as an object derived from the \texttt{PlusResult<T>} class:

\begin{itemize}
    \item \texttt{PlusSuccess<T>} -- successful operation:
    \begin{itemize}
        \item \texttt{data} -- operation result of type \texttt{T}
        \item \texttt{isSuccess} -- always \texttt{true}
        \item \texttt{errorCode} -- always \texttt{null}
        \item \texttt{message} -- always \texttt{null}
        \item \texttt{exception} -- always \texttt{null}
    \end{itemize}
    \item \texttt{PlusError<T>} -- operation error:
    \begin{itemize}
        \item \texttt{data} -- always \texttt{null}
        \item \texttt{isSuccess} -- always \texttt{false}
        \item \texttt{errorCode} -- error code of type \texttt{PlusErrorCode}
        \item \texttt{message} -- optional error description (\texttt{String?})
        \item \texttt{exception} -- optional exception (\texttt{Object?})
    \end{itemize}
\end{itemize}

\textbf{Common properties:}
\begin{itemize}
    \item \texttt{isSuccess} -- \texttt{bool} determining whether the operation was successful
    \item \texttt{data} -- operation result or \texttt{null} in case of error
    \item \texttt{errorCode} -- error code or \texttt{null} in case of success
    \item \texttt{message} -- error description or \texttt{null} in case of success
    \item \texttt{exception} -- exception or \texttt{null} in case of success
\end{itemize}

\newpage

\subsection{Types of Operation Results}
The NFC Services system uses the \textbf{result handling pattern} for consistent management of operation successes and errors:

\subsubsection*{Result Structure}
\begin{lstlisting}[language=Java]
// Success
PlusSuccess<String>(
  data: "48656C6C6F20576F726C642100000000"
)

// Error
PlusError<String>(
  error: PlusErrorCode.authenticationFailed,
  message: "Authentication failed - check your key"
)
\end{lstlisting}

\subsubsection*{Checking Result Type}
\begin{lstlisting}[language=Java]
final result = await plus.readPlus(...);

if (result.isSuccess) {
    final data = result.data!;
    print('Read data: $data');
} else {
    final error = result.errorCode!;
    print('Error: ${error.description}');
}

if (!result.isSuccess) {
    final error = result.errorCode!;
    print('Error: ${error.description}');
} else {
    final data = result.data!;
    print('Read data: $data');
}
\end{lstlisting}

\subsubsection*{Properties Available for All Results}
\begin{itemize}
    \item \texttt{isSuccess} -- \texttt{bool} - whether the operation succeeded
    \item \texttt{data} -- \texttt{T?} - result data (null for errors)
    \item \texttt{errorCode} -- \texttt{PlusErrorCode?} - error code (null for successes)
\end{itemize}

\newpage

\subsection{Constructor}
\begin{lstlisting}[language=Java]
PlusServices({
  bool debugMode = false, 
  int timeout = 10, 
  String alertMessage, 
  String multipleTagMessage
})
\end{lstlisting}

\textbf{Parameters:}
\begin{itemize}
    \item \texttt{debugMode} (\texttt{bool}) -- enables detailed logging
    \item \texttt{timeout} (\texttt{int}) -- NFC operation timeout (seconds)
    \item \texttt{alertMessage} (\texttt{String}) -- message during scanning (iOS only)
    \item \texttt{multipleTagMessage} (\texttt{String}) -- message for multiple tags (iOS only)
\end{itemize}

\subsection{Public Methods}
\begin{lstlisting}[language=Java]
Future<PlusResult<String>> readPlus(
  int sector,
  int block,
  String keyType,
  String keyHex, {
  bool isPlaintext = true,
  bool requiresMacOnResponse = false,
})
\end{lstlisting}

\textbf{Parameters:}
\begin{itemize}
    \item \texttt{sector} (\texttt{int}) -- sector number (1-39)
    \item \texttt{block} (\texttt{int}) -- block number (0-14)
    \item \texttt{keyType} (\texttt{String}) -- key type ("Key A" or "Key B")
    \item \texttt{keyHex} (\texttt{String}) -- AES key in hex format (32 characters)
    \item \texttt{isPlaintext} (\texttt{bool}) -- whether to read in plaintext mode
    \item \texttt{requiresMacOnResponse} (\texttt{bool}) -- whether MAC is required in response
\end{itemize}

\subsection{Helper Types}
\begin{itemize}
    \item \texttt{PlusAuthResult} -- authentication result (codes, random numbers, transaction ID)
    \item \texttt{PlusSessionKeys} -- session keys (kenc, kmac)
    \item \texttt{PlusReadResult} -- block read result (code, data, description)
    \item \texttt{KeyType} -- key type (keyA, keyB)
\end{itemize}

\subsection{Private Methods}

\subsubsection*{Validation and Error Handling}
\begin{itemize}
    \item \texttt{PlusErrorCode \_validatePlusReadParameters(...)}\\
    Validates input parameter correctness (sector, block, key type, hex format)
    
    \item \texttt{String \_getNfcErrorMessage(String code)}\\
    Maps NFC error code to readable text description
    
    \item \texttt{void \_log(String message, \{String level = 'INFO'\})}\\
    Logs messages (active only in \texttt{debugMode})
\end{itemize}

\subsubsection*{Authentication}
\begin{itemize}
    \item \texttt{Future<PlusResult<Uint8List>> \_processFirstAuthStep(...)}\\
    First authentication step (getting and decrypting RandomB)
    
    \item \texttt{Future<PlusResult<PlusAuthResult>> \_processSecondAuthStep(...)}\\
    Second authentication step (RandomA, session, verification)
    
    \item \texttt{Future<PlusResult<PlusAuthResult>> \_authenticateSector(...)}\\
    Complete sector authentication with retry handling
    
    \item \texttt{PlusResult<PlusAuthResult> \_validateAuthResult(...)}\\
    Authentication result validation and transaction ID extraction
\end{itemize}

\subsubsection*{Session Key Generation}
\begin{itemize}
    \item \texttt{PlusResult<PlusSessionKeys> \_generateSessionKeys(...)}\\
    Generates session keys (kenc, kmac) based on authentication data
\end{itemize}

\subsubsection*{Block Operations}
\begin{itemize}
    \item \texttt{Future<PlusResult<PlusReadResult>> \_readBlockData(...)}\\
    Reads block data with CMAC support
    
    \item \texttt{String \_decryptBlockData(...)}\\
    Decrypts block data using session keys
    
    \item \texttt{String \_extractBlockData(...)}\\
    Extracts specific block data from full response
\end{itemize}

\subsubsection*{Command Generation}
\begin{itemize}
    \item \texttt{String \_generateAuthCommand(...)}\\
    Generates authentication command (70h)
    
    \item \texttt{String \_generateBlockReadCommandCode(...)}\\
    Generates block read command code (31h, 32h, 33h)
    
    \item \texttt{String \_generateBlockReadCommand(...)}\\
    Generates complete block read command
    
    \item \texttt{String \_calculateCmac(...)}\\
    Calculates CMAC for command authentication
\end{itemize}

\subsection{Error Handling}
Each method returns a detailed error code (\texttt{PlusErrorCode}) and optional description and exception. Error codes are mapped to readable messages, making it easier to diagnose and handle exceptions in the application.

\begin{table}[h!]
\centering
\begin{tabular}{|l|l|p{6cm}|}
\hline
\textbf{Error Code} & \textbf{Code} & \textbf{Description} \\
\hline
\texttt{success} & 00 & Operation completed successfully \\
\hline
\texttt{invalidSector} & PPS & Invalid sector number (must be 1-39) \\
\hline
\texttt{invalidBlock} & PPB & Invalid block number (must be 0-14) \\
\hline
\texttt{invalidKeyType} & PPKT & Invalid key type \\
\hline
\texttt{invalidKeyHex} & PPK & Invalid AES key hex format (must be 32 hex characters) \\
\hline
\texttt{timeoutError} & PPT & Operation timeout exceeded \\
\hline
\texttt{nfcUnavailable} & PPN & NFC not available \\
\hline
\texttt{sessionError} & PPSN & NFC session error \\
\hline
\texttt{unexpectedState} & PPU & Unexpected application state \\
\hline
\texttt{userCancelled} & PPUC & User cancelled operation \\
\hline
\texttt{cardError} & CARD & Card error \\
\hline
\texttt{authenticationFailed} & PPAF & Authentication error \\
\hline
\texttt{sessionKeyError} & PPSE & Session key generation error \\
\hline
\texttt{cmacError} & PPCM & CMAC calculation error \\
\hline
\texttt{dataDecryptionError} & PPDD & Data decryption error \\
\hline
\texttt{invalidResponse} & PPIR & Invalid card response \\
\hline
\texttt{cardConnectionLost} & PCCL & Card connection lost - keep card steady \\
\hline
\texttt{communicationError} & PCE & Card communication error \\
\hline
\texttt{unknownError} & UE & Unknown error. Try again \\
\hline
\end{tabular}
\caption{\texttt{PlusErrorCode} error codes and their meanings}
\end{table}


\newpage

\subsection{Usage Example}
\begin{lstlisting}[language=Java]
final plus = PlusServices(debugMode: true);
final result = await plus.readPlus(
  1,
  0,
  "Key A",
  "0123456789ABCDEF0123456789ABCDEF",
  isPlaintext: true,
  requiresMacOnResponse: false,
);

if (result.isSuccess) {
  print('Block data: ${result.data}');
  print('Operation successful: ${result.isSuccess}');
} else {
  print('Error code: ${result.errorCode}');
  print('Error message: ${result.message}');
  if (result.exception != null) {
    print('Exception: ${result.exception}');
  }
}
\end{lstlisting}

\subsection{What readPlus Can Return}

The \texttt{readPlus} method returns an object of type \texttt{PlusResult<String>}, which can be one of two types:

\subsubsection*{PlusSuccess<String> - Successful Operation}
\begin{itemize}
    \item \texttt{isSuccess} = \texttt{true}
    \item \texttt{data} = \texttt{String} - data read from block in hex format
    \item \texttt{errorCode} = \texttt{null}
    \item \texttt{message} = \texttt{null}
    \item \texttt{exception} = \texttt{null}
\end{itemize}

\textbf{Success example:}
\begin{lstlisting}[language=Java]
PlusSuccess<String>(
  data: "48656C6C6F20576F726C642100000000" // "Hello World!" in hex
)
\end{lstlisting}

\newpage

\subsubsection*{PlusError<String> - Operation Error}
\begin{itemize}
    \item \texttt{isSuccess} = \texttt{false}
    \item \texttt{data} = \texttt{null}
    \item \texttt{errorCode} = \texttt{PlusErrorCode} - error code
    \item \texttt{message} = \texttt{String?} - optional error description
    \item \texttt{exception} = \texttt{Object?} - optional exception
\end{itemize}

\textbf{Error examples:}
\begin{lstlisting}[language=Java]
// Error: parameter validation
PlusError<String>(
  error: PlusErrorCode.invalidSector,
  message: "Invalid sector (must be 1-39)"
)

// Error: NFC not available
PlusError<String>(
  error: PlusErrorCode.nfcUnavailable,
  message: "Please enable NFC in your device settings"
)

// Error: authentication failed
PlusError<String>(
  error: PlusErrorCode.authenticationFailed,
  message: "Authentication failed - check your key"
)

// Error: connection lost
PlusError<String>(
  error: PlusErrorCode.cardConnectionLost,
  message: "Card connection lost - please keep card steady"
)

// Error: unknown error with exception
PlusError<String>(
  error: PlusErrorCode.unknownError,
  message: "Unknown error. Try again",
  exception: TimeoutException("Operation timed out")
)
\end{lstlisting}

\newpage

\subsubsection*{Returned Data Format}
In case of success, \texttt{data} contains:
\begin{itemize}
    \item \textbf{Format:} String in hex format (32 characters = 16 bytes)
    \item \textbf{Example:} \texttt{"48656C6C6F20576F726C642100000000"}
    \item \textbf{Decryption:} Data is automatically decrypted before returning
\end{itemize}

\subsubsection*{Error Handling in Application}
\begin{lstlisting}[language=Java]
final result = await plus.readPlus(sector, block, keyType, keyHex);

if (result.isSuccess) {
    // Success - use result.data
    final blockData = result.data!;
    print('Read data: $blockData');
} else {
    // Error - check code and message
    final errorCode = result.errorCode!;
    final errorMessage = result.message;
    
    switch (errorCode) {
        case PlusErrorCode.nfcUnavailable:
            showDialog('Enable NFC in device settings');
            break;
        case PlusErrorCode.cardConnectionLost:
            showDialog('Keep the card steady near the phone');
            break;
        case PlusErrorCode.authenticationFailed:
            showDialog('Invalid authentication key');
            break;
        default:
            showDialog('Error: ${errorMessage ?? errorCode.description}');
    }
}
\end{lstlisting}

\section{Shared Modules}

The NFC Services system contains shared modules that provide common functionality for both services (MIFARE DESFire and MIFARE Plus). These modules contain cryptographic tools, binary conversions, and basic NFC operations.

\subsection{binary\_conversion.dart Module}

\subsubsection{Class Description}
\texttt{BinaryConversion} is a utility class for converting data between different formats. It provides methods for converting hex strings to binary, hex, and ASCII formats with support for byte range selection and byte reversal.

\subsubsection{Main Functionality}
\begin{itemize}
    \item \textbf{Binary format conversion} -- \texttt{toBinary()} converts hex to binary string with spaces between bytes
    \item \textbf{Hex format conversion} -- \texttt{toHex()} returns hex with spaces between bytes in uppercase format
    \item \textbf{ASCII conversion} -- \texttt{toASCII()} converts hex to readable ASCII text with invalid characters replaced by dots
\end{itemize}

\subsubsection{Usage Examples}
\begin{lstlisting}[language=Java]
// Binary conversion
final binary = BinaryConversion.toBinary("12345678", 0, 1);
// Result: "00010010 00110100"

// Hex conversion
final hex = BinaryConversion.toHex("12345678", 0, 1);
// Result: "12 34"

// ASCII conversion
final ascii = BinaryConversion.toASCII("48656C6C6F", 0, 4);
// Result: "Hello"
\end{lstlisting}

\subsubsection{Error Handling}
The class includes detailed input parameter validation:
\begin{itemize}
    \item Checking if string is not empty
    \item Verifying even length of hex string
    \item Validating correct hex characters
    \item Controlling byte position range
\end{itemize}

\subsection{crypto\_utilities.dart Module}

\subsubsection{Class Description}
\texttt{CryptoUtilities} is a utility class for cryptographic operations. It provides methods for AES encryption, random number generation, byte manipulation, and CMAC calculation for MIFARE Plus operations.

\subsubsection{Main Functionality}

\paragraph{AES Encryption}
\begin{itemize}
    \item \texttt{aesEncrypt()} -- encrypts hex data with key and IV in AES-CBC mode
    \item \texttt{aesDecrypt()} -- decrypts hex data with key and IV in AES-CBC mode
    \item CBC mode support without padding for NFC operations
    \item Parameter validation (32 hex characters for key and IV)
\end{itemize}

\paragraph{Random Number Generation}
\begin{itemize}
    \item \texttt{generateSecureRandom()} -- generates cryptographically secure random bytes
    \item Uses \texttt{Random.secure()} for cryptographic operations
    \item Typically used for generating RandomA and RandomB in NFC protocols
\end{itemize}

\paragraph{Byte Manipulation}
\begin{itemize}
    \item \texttt{rotateLeft()} -- rotates byte array one position to the left
    \item \texttt{rotateRight()} -- rotates byte array one position to the right
    \item \texttt{xorBytes()} -- performs XOR operation on two byte arrays
    \item \texttt{padToBlockSize()} -- adds PKCS\#7 padding to data
    \item \texttt{removePadding()} -- removes PKCS\#7 padding from data
\end{itemize}

\paragraph{Numeric Conversions}
\begin{itemize}
    \item \texttt{intToLittleEndianHex()} -- converts number to 4-character little-endian hex
    \item \texttt{littleEndianHexToInt()} -- converts little-endian hex to number
    \item Used for formatting NFC commands where little-endian is required
\end{itemize}

\paragraph{CMAC Calculation}
\begin{itemize}
    \item \texttt{computeAesCmac()} -- computes AES-CMAC with specified output size
    \item Supports 128, 192, and 256-bit sizes
    \item Used for integrity verification in MIFARE Plus operations
\end{itemize}

\subsubsection{Usage Examples}
\begin{lstlisting}[language=Java]
// AES encryption
final encrypted = CryptoUtilities.aesEncrypt(
    "48656C6C6F20576F726C64", // "Hello World" in hex
    "0123456789ABCDEF0123456789ABCDEF", // 32-char key
    "00000000000000000000000000000000"  // 32-char IV
);

// Random byte generation
final randomBytes = CryptoUtilities.generateSecureRandom(16);

// CMAC calculation
final cmac = CryptoUtilities.computeAesCmac(
    "0123456789ABCDEF0123456789ABCDEF", // key
    "48656C6C6F20576F726C64", // data
    128 // output bits
);

// Little-endian conversion
final leHex = CryptoUtilities.intToLittleEndianHex(258);
// Result: "0201" (not "0102")

// Byte XOR
final result = CryptoUtilities.xorBytes(
    Uint8List.fromList([1, 2, 3]),
    Uint8List.fromList([4, 5, 6])
);
// Result: [5, 7, 5]
\end{lstlisting}

\subsubsection{Error Handling}
The class includes its own exception type \texttt{CryptoException} and detailed validation:
\begin{itemize}
    \item Checking correctness of AES parameters (length, hex format)
    \item CMAC parameter validation (output size, data format)
    \item Byte array size control for XOR operations
    \item PKCS\#7 padding verification
\end{itemize}

\subsubsection{External Dependencies}
The module uses the following libraries:
\begin{itemize}
    \item \texttt{convert} -- hex/binary conversions
    \item \texttt{encrypt} -- AES cryptographic operations
    \item \texttt{pointycastle} -- CMAC calculation and advanced cryptographic functions
\end{itemize}

\subsection{nfc\_utilities.dart Module}

\subsubsection{Class Description}
\texttt{nfc\_utilities.dart} is a module containing basic tools for handling NFC operations. It provides functions for checking NFC availability, card detection, command sending, and communication error handling with cards.

\subsubsection{Main Functionality}
\begin{itemize}
    \item \textbf{NFC availability checking} -- \texttt{isNfcAvailable()} verifies if NFC is available, disabled, or unsupported
    \item \textbf{Card detection} -- \texttt{pollForCard()} scans and detects NFC cards with timeout support
    \item \textbf{Command sending} -- \texttt{sendCommand()} and \texttt{sendCommands()} send hex commands to card
    \item \textbf{Response validation} -- \texttt{isSuccessfulResponse()} checks card response correctness
    \item \textbf{Error handling} -- \texttt{NfcErrorCode} and \texttt{NfcResult} provide consistent error handling
    \item \textbf{Debug logging} -- \texttt{debugPrint()} and \texttt{debugPrintWithLevel()} for diagnostics
\end{itemize}

\subsubsection{Supported Card Types}
The module supports the following NFC card types:
\begin{itemize}
    \item \texttt{mifarePlus} -- MIFARE Plus SL3 cards
    \item \texttt{mifareDesfire} -- MIFARE DESFire cards
\end{itemize}

\subsubsection{Usage Example}
\begin{lstlisting}[language=Java]
// Check NFC availability
final nfcResult = await isNfcAvailable();
if (!nfcResult.isSuccess) {
    print('NFC error: ${nfcResult.message}');
    return;
}

// Detect card
final cardResult = await pollForCard(
    debugMode: true,
    timeout: 10,
    alertMessage: "Please place card near device",
    multipleTagMessage: "Multiple cards detected"
);
\end{lstlisting}

\end{document} 